\documentclass[12pt]{article}
\usepackage[utf8]{inputenc}
\usepackage[T1]{fontenc}
\usepackage{hyperref}
\title{Tz'utujil Language Resources for K-12 Curriculum Development}
\author{Compiled by Claude Sonnet 5 (claude-sonnet-5) for Andr\'as Kornai}
\date{}
\begin{document}
\maketitle
\section{Introduction}
Tz'utujil (ISO 639-3 tzj) is a Mayan language of Guatemala with about 60,000 speakers. Its regulator is the Academia de Lenguas Mayas de Guatemala (ALMG).
\section{Primary Educational Resources}
No primers or textbooks were found in this search.
\section{Linguistic and Reference Materials}
\begin{itemize}
\item Vocabulario Tz'utujil (ALMG, 2025), a bilingual Tz'utujil--Spanish vocabulary. \url{https://almg.org.gt/wp-content/uploads/2025/09/VOCABULARIO-TZUTUJIL-VERSION-FINAL-2025.pdf}
\item Glottolog's language page, whose bibliography lists Dayley's \emph{Tzutujil Grammar} (1981, 1985), Garc\'ia Ixmat\'a's \emph{Gram\'atica Tz'utujiil} (1997), and the \emph{Diccionario Tz'utujil} (1996). \url{https://glottolog.org/resource/languoid/iso/tzj}
\item ASJP word list for Santiago Atitl\'an Tzutujil. \url{https://asjp.clld.org/languages/TZUTUJIL_SANTIAGO_ATITLAN}
\end{itemize}
\section{Religious and Community Resources}
None found.
\section{Digital and Audio Resources}
\begin{itemize}
\item Endangered Languages Project entry. \url{https://endangeredlanguages.com/node/97301}
\item Overview of the language. \url{https://wikipedia.com/wiki/ISO_639:tzj}
\end{itemize}
\section{Curriculum Development Context}
The ALMG's standard orthography and its 2025 vocabulary are the natural starting points. Grammars and a dictionary exist, but I did not see online copies. No school-level materials were found in this search.
\end{document}